% ============================================================
%  课程笔记封面模板 —— 「富春山居图」风格
%  取材：黄公望《富春山居图》（元·水墨长卷）：米黄绢底、披麻皴的
%        淡墨远山、留白的江水、疏落的秋树与一叶扁舟、朱红钤印。
%  与 cover_Impression风.tex 同构（眉题 / 主标题 / 菱形饰线 /
%  副标题 / 底部作者日期），直接 \input{} 复用：
%      % ============================================================
%  课程笔记封面模板 —— 「富春山居图」风格
%  取材：黄公望《富春山居图》（元·水墨长卷）：米黄绢底、披麻皴的
%        淡墨远山、留白的江水、疏落的秋树与一叶扁舟、朱红钤印。
%  与 cover_Impression风.tex 同构（眉题 / 主标题 / 菱形饰线 /
%  副标题 / 底部作者日期），直接 \input{} 复用：
%      % ============================================================
%  课程笔记封面模板 —— 「富春山居图」风格
%  取材：黄公望《富春山居图》（元·水墨长卷）：米黄绢底、披麻皴的
%        淡墨远山、留白的江水、疏落的秋树与一叶扁舟、朱红钤印。
%  与 cover_Impression风.tex 同构（眉题 / 主标题 / 菱形饰线 /
%  副标题 / 底部作者日期），直接 \input{} 复用：
%      % ============================================================
%  课程笔记封面模板 —— 「富春山居图」风格
%  取材：黄公望《富春山居图》（元·水墨长卷）：米黄绢底、披麻皴的
%        淡墨远山、留白的江水、疏落的秋树与一叶扁舟、朱红钤印。
%  与 cover_Impression风.tex 同构（眉题 / 主标题 / 菱形饰线 /
%  副标题 / 底部作者日期），直接 \input{} 复用：
%      \input{cover_富春山居图风}
%
%  配色：fc* 前缀（绢素 + 水墨五色）。本文件用 \providecolor 兜底，
%        单独复制到任何笔记本也能编译；想整体换调改这里即可。
%
%  注意：整幅画面绘于 \begin{scope}[shift={(current page.south west)}] 内，
%        所有坐标以「页面左下角为原点、单位 cm」书写；脱离该 scope
%        直接写 (x,y)，图形会落到页面外而不可见。
%
%  可用字段（在主文件导言区用 \newcommand 预定义即可覆盖缺省值）：
%      \notename     主标题（必填，主模板已定义）
%      \noteauthor   作者（必填，主模板已定义）
%      \notecourseen 眉题英文（缺省 Course Notes）
%      \notesubtitle 副标题（缺省「学习笔记」）
%      \noteextra    底部附加信息（缺省不显示）
% ============================================================

% ---- 《富春山居图》取色（绢素 + 水墨） ----
\providecolor{fcSilk}{RGB}{236,226,202}     % 绢底：米黄
\providecolor{fcSilkDark}{RGB}{214,202,174} % 绢底：暗部
\providecolor{fcInk3}{RGB}{178,176,168}     % 淡墨（远山）
\providecolor{fcInk2}{RGB}{116,114,108}     % 中墨（近坡）
\providecolor{fcInk}{RGB}{54,52,48}         % 浓墨（树、点、标题）
\providecolor{fcSeal}{RGB}{168,52,44}       % 朱红钤印

% ---- 缺省字段（主文件已定义的同名命令不会被覆盖） ----
\providecommand{\notecourseen}{Course Notes}
\providecommand{\notesubtitle}{学\ 习\ 笔\ 记}
\providecommand{\noteextra}{}

\begin{titlepage}
	\begin{tikzpicture}[remember picture, overlay]
		\begin{scope}[shift={(current page.south west)}]

		% ============================================================
		%  1. 绢底（米黄，下缘略深；带细微斑驳）
		% ============================================================
		\shade[top color=fcSilk, bottom color=fcSilkDark] (0,0) rectangle (21,29.7);
		\foreach \x/\y in {3.2/26.4, 8.6/24.2, 13.4/27.0, 17.8/25.2, 6.0/21.0,
		                   11.2/19.4, 19.0/21.8, 2.0/13.6, 15.2/8.2}{
			\fill[fcInk3, opacity=0.10] (\x,\y) ellipse[x radius=0.9, y radius=0.35];
		}

		% ============================================================
		%  2. 远山（三层，越远越淡；披麻皴短线）
		% ============================================================
		\fill[fcInk3, opacity=0.40]
			(0,20.0) -- (2.4,22.6) -- (4.0,21.2) -- (6.2,23.8) -- (8.6,21.6)
			-- (11.0,24.2) -- (13.4,21.4) -- (15.8,23.4) -- (18.2,21.0)
			-- (21,22.4) -- (21,17.4) -- (0,17.4) -- cycle;
		\fill[fcInk3, opacity=0.62]
			(0,17.8) -- (2.8,20.0) -- (4.6,18.6) -- (7.0,20.8) -- (9.4,18.4)
			-- (12.2,20.6) -- (15.0,18.2) -- (18.0,20.0) -- (21,18.2)
			-- (21,15.0) -- (0,15.0) -- cycle;
		\fill[fcInk2, opacity=0.45]
			(0,15.4) -- (2.2,17.0) -- (4.0,15.8) -- (6.4,17.4) -- (8.8,15.6)
			-- (11.6,17.2) -- (14.4,15.4) -- (17.6,16.8) -- (21,15.0)
			-- (21,12.6) -- (0,12.6) -- cycle;
		% 披麻皴：山体上的短竖笔
		\foreach \x/\y/\l in {3.0/19.4/1.5, 5.4/20.2/1.8, 7.8/19.6/1.6, 10.4/20.4/2.0,
		                      13.0/19.0/1.7, 16.0/19.4/1.5, 18.6/18.6/1.4}{
			\draw[fcInk2, opacity=0.35, line width=0.9pt]
				(\x,\y) .. controls (\x+0.18,\y-\l*0.45) and (\x-0.12,\y-\l*0.75) .. (\x+0.10,\y-\l);
		}

		% ============================================================
		%  3. 江水（留白 + 数笔淡墨波纹）
		% ============================================================
		\foreach \y in {12.0, 11.2, 10.4}{
			\draw[fcInk3, opacity=0.42, line width=0.8pt]
				plot[domain=1.2:19.6, samples=40, smooth]
					(\x, {\y + 0.13*sin(deg(2.0*\x + 3*\y))});
		}

		% ============================================================
		%  4. 一叶扁舟
		% ============================================================
		\fill[fcInk]
			(13.4,11.15) .. controls (14.0,10.90) and (15.0,10.90) .. (15.7,11.15)
			.. controls (15.2,11.42) and (13.9,11.42) .. (13.4,11.15) -- cycle;
		\draw[fcInk, line width=0.9pt] (14.6,11.15) -- (14.6,11.95);
		\fill[fcInk] (14.6,12.15) circle[radius=0.13];

		% ============================================================
		%  5. 近景坡石（左右两岸，中墨勾勒 + 皴笔）
		% ============================================================
		\fill[fcInk2, opacity=0.75]
			(0,0) -- (0,7.2) -- (1.6,8.6) -- (3.2,7.8) -- (4.8,9.4)
			-- (6.4,8.2) -- (8.0,9.0) -- (9.2,7.4) -- (10.2,5.6)
			-- (11.0,3.2) -- (11.6,0) -- cycle;
		\fill[fcInk2, opacity=0.70]
			(21,0) -- (21,6.0) -- (19.4,7.4) -- (17.8,6.6) -- (16.2,8.0)
			-- (14.6,6.8) -- (13.2,5.2) -- (12.4,2.6) -- (12.0,0) -- cycle;
		\draw[fcInk, opacity=0.55, line width=1.1pt]
			(0,7.2) -- (1.6,8.6) -- (3.2,7.8) -- (4.8,9.4) -- (6.4,8.2)
			-- (8.0,9.0) -- (9.2,7.4) -- (10.2,5.6) -- (11.0,3.2);
		\draw[fcInk, opacity=0.55, line width=1.1pt]
			(21,6.0) -- (19.4,7.4) -- (17.8,6.6) -- (16.2,8.0) -- (14.6,6.8)
			-- (13.2,5.2) -- (12.4,2.6);
		\foreach \x/\y/\l in {2.0/7.6/1.4, 4.2/8.2/1.6, 6.6/7.4/1.3, 8.6/7.8/1.2,
		                      15.0/6.2/1.3, 17.0/6.8/1.5, 19.0/6.2/1.2}{
			\draw[fcInk, opacity=0.30, line width=0.8pt] (\x,\y) -- (\x+0.12,\y-\l);
		}

		% ============================================================
		%  6. 疏落的秋树（竖笔为干，横点为叶）
		% ============================================================
		\foreach \x/\h in {1.4/3.2, 2.2/4.3, 3.1/3.6, 3.9/2.8, 15.6/3.0, 16.5/4.0, 17.3/3.2}{
			\draw[fcInk, line width=1.4pt] (\x,8.6) -- (\x,{8.6+\h});
			\foreach \dx/\dy in {-0.42/0.30, -0.16/0.52, 0.16/0.44, 0.42/0.24}{
				\fill[fcInk, opacity=0.75]
					({\x+\dx},{8.6+\h+\dy}) ellipse[x radius=0.20, y radius=0.10];
			}
		}

		% ============================================================
		%  7. 标题衬底（一层薄绢，保证文字可读）
		% ============================================================
		\fill[fcSilk, opacity=0.50]
			(2.6,16.9) -- (17.4,17.3) -- (18.0,22.6) -- (16.4,23.6) -- (2.2,22.9) -- cycle;

		% ============================================================
		%  8. 朱红钤印（右上、左下各一）
		% ============================================================
		\fill[fcSeal, opacity=0.85] (18.3,25.0) rectangle (19.7,26.4);
		\draw[fcSilk, line width=0.7pt, opacity=0.9] (18.55,25.25) rectangle (19.45,26.15);
		\fill[fcSeal, opacity=0.80] (1.5,3.0) rectangle (2.7,4.2);
		\draw[fcSilk, line width=0.6pt, opacity=0.9] (1.72,3.22) rectangle (2.48,3.98);

		% ============================================================
		%  9. 细边框（墨线）
		% ============================================================
		\draw[fcInk, line width=0.7pt, opacity=0.45]
			([shift={(0.85,0.85)}]current page.south west) rectangle
			([shift={(-0.85,-0.85)}]current page.north east);

		% ============================================================
		% 10. 顶部眉题（课程英文小字 + 自适应墨线）
		% ============================================================
		\node[anchor=north, text=fcInk, align=center]
			(fcEyebrow) at ([yshift=-3.4cm]current page.north) {%
			{\sffamily\fontsize{12}{15}\selectfont \uppercase{\notecourseen}}
		};
		\draw[fcInk, line width=1.0pt] ($(fcEyebrow.west)+(-0.25cm,0)$) -- ++(-1.1cm,0);
		\draw[fcInk, line width=1.0pt] ($(fcEyebrow.east)+(0.25cm,0)$) -- ++(1.1cm,0);

		% ============================================================
		% 11. 主标题（居中，使用 \notename）
		% ============================================================
		\node[anchor=center, align=center, text=fcInk]
			at ([yshift=-8.3cm]current page.north) {%
			{\fontsize{38}{48}\sffamily\bfseries \notename}
		};

		% ============================================================
		% 12. 菱形 + 双细线
		% ============================================================
		\coordinate (C) at ([yshift=-11.6cm]current page.north);
		\draw[fcInk, line width=1.1pt] (C) ++(-2.2cm,0) -- ++(1.55cm,0);
		\draw[fcInk, line width=1.1pt] (C) ++(0.65cm,0) -- ++(1.55cm,0);
		\node[fill=fcInk, minimum size=0.24cm, inner sep=0pt, rotate=45] at (C) {};

		% ============================================================
		% 13. 副标题（使用 \notesubtitle）
		% ============================================================
		\node[anchor=north, text=fcInk, align=center]
			at ([yshift=-12.4cm]current page.north) {
			{\fontsize{15}{20}\sffamily \notesubtitle}
		};

		% ============================================================
		% 14. 底部作者、日期、附加信息
		% ============================================================
		\coordinate (B) at ([yshift=2.7cm]current page.south);
		\draw[fcInk, line width=0.8pt, opacity=0.8] (B) ++(-1.1cm,0.95cm) -- ++(2.2cm,0);
		\node[anchor=south, text=fcInk, align=center] at (B) {
			\begin{tabular}{c}
				{\fontsize{19}{24}\sffamily\bfseries \noteauthor} \\[0.35cm]
				{\large \sffamily \today}
				\ifstrempty{\noteextra}{}{\\[0.3cm]{\normalsize\sffamily \noteextra}}
			\end{tabular}
		};

		% ============================================================
		% 15. 右下角落款（画作信息）
		% ============================================================
		\node[anchor=south east, text=fcInk, opacity=0.65, font=\footnotesize\itshape]
			at ([shift={(-1.35,1.25)}]current page.south east)
			{Huang Gongwang, Dwelling in the Fuchun Mountains, 1350};

		\end{scope}
	\end{tikzpicture}
\end{titlepage}

%
%  配色：fc* 前缀（绢素 + 水墨五色）。本文件用 \providecolor 兜底，
%        单独复制到任何笔记本也能编译；想整体换调改这里即可。
%
%  注意：整幅画面绘于 \begin{scope}[shift={(current page.south west)}] 内，
%        所有坐标以「页面左下角为原点、单位 cm」书写；脱离该 scope
%        直接写 (x,y)，图形会落到页面外而不可见。
%
%  可用字段（在主文件导言区用 \newcommand 预定义即可覆盖缺省值）：
%      \notename     主标题（必填，主模板已定义）
%      \noteauthor   作者（必填，主模板已定义）
%      \notecourseen 眉题英文（缺省 Course Notes）
%      \notesubtitle 副标题（缺省「学习笔记」）
%      \noteextra    底部附加信息（缺省不显示）
% ============================================================

% ---- 《富春山居图》取色（绢素 + 水墨） ----
\providecolor{fcSilk}{RGB}{236,226,202}     % 绢底：米黄
\providecolor{fcSilkDark}{RGB}{214,202,174} % 绢底：暗部
\providecolor{fcInk3}{RGB}{178,176,168}     % 淡墨（远山）
\providecolor{fcInk2}{RGB}{116,114,108}     % 中墨（近坡）
\providecolor{fcInk}{RGB}{54,52,48}         % 浓墨（树、点、标题）
\providecolor{fcSeal}{RGB}{168,52,44}       % 朱红钤印

% ---- 缺省字段（主文件已定义的同名命令不会被覆盖） ----
\providecommand{\notecourseen}{Course Notes}
\providecommand{\notesubtitle}{学\ 习\ 笔\ 记}
\providecommand{\noteextra}{}

\begin{titlepage}
	\begin{tikzpicture}[remember picture, overlay]
		\begin{scope}[shift={(current page.south west)}]

		% ============================================================
		%  1. 绢底（米黄，下缘略深；带细微斑驳）
		% ============================================================
		\shade[top color=fcSilk, bottom color=fcSilkDark] (0,0) rectangle (21,29.7);
		\foreach \x/\y in {3.2/26.4, 8.6/24.2, 13.4/27.0, 17.8/25.2, 6.0/21.0,
		                   11.2/19.4, 19.0/21.8, 2.0/13.6, 15.2/8.2}{
			\fill[fcInk3, opacity=0.10] (\x,\y) ellipse[x radius=0.9, y radius=0.35];
		}

		% ============================================================
		%  2. 远山（三层，越远越淡；披麻皴短线）
		% ============================================================
		\fill[fcInk3, opacity=0.40]
			(0,20.0) -- (2.4,22.6) -- (4.0,21.2) -- (6.2,23.8) -- (8.6,21.6)
			-- (11.0,24.2) -- (13.4,21.4) -- (15.8,23.4) -- (18.2,21.0)
			-- (21,22.4) -- (21,17.4) -- (0,17.4) -- cycle;
		\fill[fcInk3, opacity=0.62]
			(0,17.8) -- (2.8,20.0) -- (4.6,18.6) -- (7.0,20.8) -- (9.4,18.4)
			-- (12.2,20.6) -- (15.0,18.2) -- (18.0,20.0) -- (21,18.2)
			-- (21,15.0) -- (0,15.0) -- cycle;
		\fill[fcInk2, opacity=0.45]
			(0,15.4) -- (2.2,17.0) -- (4.0,15.8) -- (6.4,17.4) -- (8.8,15.6)
			-- (11.6,17.2) -- (14.4,15.4) -- (17.6,16.8) -- (21,15.0)
			-- (21,12.6) -- (0,12.6) -- cycle;
		% 披麻皴：山体上的短竖笔
		\foreach \x/\y/\l in {3.0/19.4/1.5, 5.4/20.2/1.8, 7.8/19.6/1.6, 10.4/20.4/2.0,
		                      13.0/19.0/1.7, 16.0/19.4/1.5, 18.6/18.6/1.4}{
			\draw[fcInk2, opacity=0.35, line width=0.9pt]
				(\x,\y) .. controls (\x+0.18,\y-\l*0.45) and (\x-0.12,\y-\l*0.75) .. (\x+0.10,\y-\l);
		}

		% ============================================================
		%  3. 江水（留白 + 数笔淡墨波纹）
		% ============================================================
		\foreach \y in {12.0, 11.2, 10.4}{
			\draw[fcInk3, opacity=0.42, line width=0.8pt]
				plot[domain=1.2:19.6, samples=40, smooth]
					(\x, {\y + 0.13*sin(deg(2.0*\x + 3*\y))});
		}

		% ============================================================
		%  4. 一叶扁舟
		% ============================================================
		\fill[fcInk]
			(13.4,11.15) .. controls (14.0,10.90) and (15.0,10.90) .. (15.7,11.15)
			.. controls (15.2,11.42) and (13.9,11.42) .. (13.4,11.15) -- cycle;
		\draw[fcInk, line width=0.9pt] (14.6,11.15) -- (14.6,11.95);
		\fill[fcInk] (14.6,12.15) circle[radius=0.13];

		% ============================================================
		%  5. 近景坡石（左右两岸，中墨勾勒 + 皴笔）
		% ============================================================
		\fill[fcInk2, opacity=0.75]
			(0,0) -- (0,7.2) -- (1.6,8.6) -- (3.2,7.8) -- (4.8,9.4)
			-- (6.4,8.2) -- (8.0,9.0) -- (9.2,7.4) -- (10.2,5.6)
			-- (11.0,3.2) -- (11.6,0) -- cycle;
		\fill[fcInk2, opacity=0.70]
			(21,0) -- (21,6.0) -- (19.4,7.4) -- (17.8,6.6) -- (16.2,8.0)
			-- (14.6,6.8) -- (13.2,5.2) -- (12.4,2.6) -- (12.0,0) -- cycle;
		\draw[fcInk, opacity=0.55, line width=1.1pt]
			(0,7.2) -- (1.6,8.6) -- (3.2,7.8) -- (4.8,9.4) -- (6.4,8.2)
			-- (8.0,9.0) -- (9.2,7.4) -- (10.2,5.6) -- (11.0,3.2);
		\draw[fcInk, opacity=0.55, line width=1.1pt]
			(21,6.0) -- (19.4,7.4) -- (17.8,6.6) -- (16.2,8.0) -- (14.6,6.8)
			-- (13.2,5.2) -- (12.4,2.6);
		\foreach \x/\y/\l in {2.0/7.6/1.4, 4.2/8.2/1.6, 6.6/7.4/1.3, 8.6/7.8/1.2,
		                      15.0/6.2/1.3, 17.0/6.8/1.5, 19.0/6.2/1.2}{
			\draw[fcInk, opacity=0.30, line width=0.8pt] (\x,\y) -- (\x+0.12,\y-\l);
		}

		% ============================================================
		%  6. 疏落的秋树（竖笔为干，横点为叶）
		% ============================================================
		\foreach \x/\h in {1.4/3.2, 2.2/4.3, 3.1/3.6, 3.9/2.8, 15.6/3.0, 16.5/4.0, 17.3/3.2}{
			\draw[fcInk, line width=1.4pt] (\x,8.6) -- (\x,{8.6+\h});
			\foreach \dx/\dy in {-0.42/0.30, -0.16/0.52, 0.16/0.44, 0.42/0.24}{
				\fill[fcInk, opacity=0.75]
					({\x+\dx},{8.6+\h+\dy}) ellipse[x radius=0.20, y radius=0.10];
			}
		}

		% ============================================================
		%  7. 标题衬底（一层薄绢，保证文字可读）
		% ============================================================
		\fill[fcSilk, opacity=0.50]
			(2.6,16.9) -- (17.4,17.3) -- (18.0,22.6) -- (16.4,23.6) -- (2.2,22.9) -- cycle;

		% ============================================================
		%  8. 朱红钤印（右上、左下各一）
		% ============================================================
		\fill[fcSeal, opacity=0.85] (18.3,25.0) rectangle (19.7,26.4);
		\draw[fcSilk, line width=0.7pt, opacity=0.9] (18.55,25.25) rectangle (19.45,26.15);
		\fill[fcSeal, opacity=0.80] (1.5,3.0) rectangle (2.7,4.2);
		\draw[fcSilk, line width=0.6pt, opacity=0.9] (1.72,3.22) rectangle (2.48,3.98);

		% ============================================================
		%  9. 细边框（墨线）
		% ============================================================
		\draw[fcInk, line width=0.7pt, opacity=0.45]
			([shift={(0.85,0.85)}]current page.south west) rectangle
			([shift={(-0.85,-0.85)}]current page.north east);

		% ============================================================
		% 10. 顶部眉题（课程英文小字 + 自适应墨线）
		% ============================================================
		\node[anchor=north, text=fcInk, align=center]
			(fcEyebrow) at ([yshift=-3.4cm]current page.north) {%
			{\sffamily\fontsize{12}{15}\selectfont \uppercase{\notecourseen}}
		};
		\draw[fcInk, line width=1.0pt] ($(fcEyebrow.west)+(-0.25cm,0)$) -- ++(-1.1cm,0);
		\draw[fcInk, line width=1.0pt] ($(fcEyebrow.east)+(0.25cm,0)$) -- ++(1.1cm,0);

		% ============================================================
		% 11. 主标题（居中，使用 \notename）
		% ============================================================
		\node[anchor=center, align=center, text=fcInk]
			at ([yshift=-8.3cm]current page.north) {%
			{\fontsize{38}{48}\sffamily\bfseries \notename}
		};

		% ============================================================
		% 12. 菱形 + 双细线
		% ============================================================
		\coordinate (C) at ([yshift=-11.6cm]current page.north);
		\draw[fcInk, line width=1.1pt] (C) ++(-2.2cm,0) -- ++(1.55cm,0);
		\draw[fcInk, line width=1.1pt] (C) ++(0.65cm,0) -- ++(1.55cm,0);
		\node[fill=fcInk, minimum size=0.24cm, inner sep=0pt, rotate=45] at (C) {};

		% ============================================================
		% 13. 副标题（使用 \notesubtitle）
		% ============================================================
		\node[anchor=north, text=fcInk, align=center]
			at ([yshift=-12.4cm]current page.north) {
			{\fontsize{15}{20}\sffamily \notesubtitle}
		};

		% ============================================================
		% 14. 底部作者、日期、附加信息
		% ============================================================
		\coordinate (B) at ([yshift=2.7cm]current page.south);
		\draw[fcInk, line width=0.8pt, opacity=0.8] (B) ++(-1.1cm,0.95cm) -- ++(2.2cm,0);
		\node[anchor=south, text=fcInk, align=center] at (B) {
			\begin{tabular}{c}
				{\fontsize{19}{24}\sffamily\bfseries \noteauthor} \\[0.35cm]
				{\large \sffamily \today}
				\ifstrempty{\noteextra}{}{\\[0.3cm]{\normalsize\sffamily \noteextra}}
			\end{tabular}
		};

		% ============================================================
		% 15. 右下角落款（画作信息）
		% ============================================================
		\node[anchor=south east, text=fcInk, opacity=0.65, font=\footnotesize\itshape]
			at ([shift={(-1.35,1.25)}]current page.south east)
			{Huang Gongwang, Dwelling in the Fuchun Mountains, 1350};

		\end{scope}
	\end{tikzpicture}
\end{titlepage}

%
%  配色：fc* 前缀（绢素 + 水墨五色）。本文件用 \providecolor 兜底，
%        单独复制到任何笔记本也能编译；想整体换调改这里即可。
%
%  注意：整幅画面绘于 \begin{scope}[shift={(current page.south west)}] 内，
%        所有坐标以「页面左下角为原点、单位 cm」书写；脱离该 scope
%        直接写 (x,y)，图形会落到页面外而不可见。
%
%  可用字段（在主文件导言区用 \newcommand 预定义即可覆盖缺省值）：
%      \notename     主标题（必填，主模板已定义）
%      \noteauthor   作者（必填，主模板已定义）
%      \notecourseen 眉题英文（缺省 Course Notes）
%      \notesubtitle 副标题（缺省「学习笔记」）
%      \noteextra    底部附加信息（缺省不显示）
% ============================================================

% ---- 《富春山居图》取色（绢素 + 水墨） ----
\providecolor{fcSilk}{RGB}{236,226,202}     % 绢底：米黄
\providecolor{fcSilkDark}{RGB}{214,202,174} % 绢底：暗部
\providecolor{fcInk3}{RGB}{178,176,168}     % 淡墨（远山）
\providecolor{fcInk2}{RGB}{116,114,108}     % 中墨（近坡）
\providecolor{fcInk}{RGB}{54,52,48}         % 浓墨（树、点、标题）
\providecolor{fcSeal}{RGB}{168,52,44}       % 朱红钤印

% ---- 缺省字段（主文件已定义的同名命令不会被覆盖） ----
\providecommand{\notecourseen}{Course Notes}
\providecommand{\notesubtitle}{学\ 习\ 笔\ 记}
\providecommand{\noteextra}{}

\begin{titlepage}
	\begin{tikzpicture}[remember picture, overlay]
		\begin{scope}[shift={(current page.south west)}]

		% ============================================================
		%  1. 绢底（米黄，下缘略深；带细微斑驳）
		% ============================================================
		\shade[top color=fcSilk, bottom color=fcSilkDark] (0,0) rectangle (21,29.7);
		\foreach \x/\y in {3.2/26.4, 8.6/24.2, 13.4/27.0, 17.8/25.2, 6.0/21.0,
		                   11.2/19.4, 19.0/21.8, 2.0/13.6, 15.2/8.2}{
			\fill[fcInk3, opacity=0.10] (\x,\y) ellipse[x radius=0.9, y radius=0.35];
		}

		% ============================================================
		%  2. 远山（三层，越远越淡；披麻皴短线）
		% ============================================================
		\fill[fcInk3, opacity=0.40]
			(0,20.0) -- (2.4,22.6) -- (4.0,21.2) -- (6.2,23.8) -- (8.6,21.6)
			-- (11.0,24.2) -- (13.4,21.4) -- (15.8,23.4) -- (18.2,21.0)
			-- (21,22.4) -- (21,17.4) -- (0,17.4) -- cycle;
		\fill[fcInk3, opacity=0.62]
			(0,17.8) -- (2.8,20.0) -- (4.6,18.6) -- (7.0,20.8) -- (9.4,18.4)
			-- (12.2,20.6) -- (15.0,18.2) -- (18.0,20.0) -- (21,18.2)
			-- (21,15.0) -- (0,15.0) -- cycle;
		\fill[fcInk2, opacity=0.45]
			(0,15.4) -- (2.2,17.0) -- (4.0,15.8) -- (6.4,17.4) -- (8.8,15.6)
			-- (11.6,17.2) -- (14.4,15.4) -- (17.6,16.8) -- (21,15.0)
			-- (21,12.6) -- (0,12.6) -- cycle;
		% 披麻皴：山体上的短竖笔
		\foreach \x/\y/\l in {3.0/19.4/1.5, 5.4/20.2/1.8, 7.8/19.6/1.6, 10.4/20.4/2.0,
		                      13.0/19.0/1.7, 16.0/19.4/1.5, 18.6/18.6/1.4}{
			\draw[fcInk2, opacity=0.35, line width=0.9pt]
				(\x,\y) .. controls (\x+0.18,\y-\l*0.45) and (\x-0.12,\y-\l*0.75) .. (\x+0.10,\y-\l);
		}

		% ============================================================
		%  3. 江水（留白 + 数笔淡墨波纹）
		% ============================================================
		\foreach \y in {12.0, 11.2, 10.4}{
			\draw[fcInk3, opacity=0.42, line width=0.8pt]
				plot[domain=1.2:19.6, samples=40, smooth]
					(\x, {\y + 0.13*sin(deg(2.0*\x + 3*\y))});
		}

		% ============================================================
		%  4. 一叶扁舟
		% ============================================================
		\fill[fcInk]
			(13.4,11.15) .. controls (14.0,10.90) and (15.0,10.90) .. (15.7,11.15)
			.. controls (15.2,11.42) and (13.9,11.42) .. (13.4,11.15) -- cycle;
		\draw[fcInk, line width=0.9pt] (14.6,11.15) -- (14.6,11.95);
		\fill[fcInk] (14.6,12.15) circle[radius=0.13];

		% ============================================================
		%  5. 近景坡石（左右两岸，中墨勾勒 + 皴笔）
		% ============================================================
		\fill[fcInk2, opacity=0.75]
			(0,0) -- (0,7.2) -- (1.6,8.6) -- (3.2,7.8) -- (4.8,9.4)
			-- (6.4,8.2) -- (8.0,9.0) -- (9.2,7.4) -- (10.2,5.6)
			-- (11.0,3.2) -- (11.6,0) -- cycle;
		\fill[fcInk2, opacity=0.70]
			(21,0) -- (21,6.0) -- (19.4,7.4) -- (17.8,6.6) -- (16.2,8.0)
			-- (14.6,6.8) -- (13.2,5.2) -- (12.4,2.6) -- (12.0,0) -- cycle;
		\draw[fcInk, opacity=0.55, line width=1.1pt]
			(0,7.2) -- (1.6,8.6) -- (3.2,7.8) -- (4.8,9.4) -- (6.4,8.2)
			-- (8.0,9.0) -- (9.2,7.4) -- (10.2,5.6) -- (11.0,3.2);
		\draw[fcInk, opacity=0.55, line width=1.1pt]
			(21,6.0) -- (19.4,7.4) -- (17.8,6.6) -- (16.2,8.0) -- (14.6,6.8)
			-- (13.2,5.2) -- (12.4,2.6);
		\foreach \x/\y/\l in {2.0/7.6/1.4, 4.2/8.2/1.6, 6.6/7.4/1.3, 8.6/7.8/1.2,
		                      15.0/6.2/1.3, 17.0/6.8/1.5, 19.0/6.2/1.2}{
			\draw[fcInk, opacity=0.30, line width=0.8pt] (\x,\y) -- (\x+0.12,\y-\l);
		}

		% ============================================================
		%  6. 疏落的秋树（竖笔为干，横点为叶）
		% ============================================================
		\foreach \x/\h in {1.4/3.2, 2.2/4.3, 3.1/3.6, 3.9/2.8, 15.6/3.0, 16.5/4.0, 17.3/3.2}{
			\draw[fcInk, line width=1.4pt] (\x,8.6) -- (\x,{8.6+\h});
			\foreach \dx/\dy in {-0.42/0.30, -0.16/0.52, 0.16/0.44, 0.42/0.24}{
				\fill[fcInk, opacity=0.75]
					({\x+\dx},{8.6+\h+\dy}) ellipse[x radius=0.20, y radius=0.10];
			}
		}

		% ============================================================
		%  7. 标题衬底（一层薄绢，保证文字可读）
		% ============================================================
		\fill[fcSilk, opacity=0.50]
			(2.6,16.9) -- (17.4,17.3) -- (18.0,22.6) -- (16.4,23.6) -- (2.2,22.9) -- cycle;

		% ============================================================
		%  8. 朱红钤印（右上、左下各一）
		% ============================================================
		\fill[fcSeal, opacity=0.85] (18.3,25.0) rectangle (19.7,26.4);
		\draw[fcSilk, line width=0.7pt, opacity=0.9] (18.55,25.25) rectangle (19.45,26.15);
		\fill[fcSeal, opacity=0.80] (1.5,3.0) rectangle (2.7,4.2);
		\draw[fcSilk, line width=0.6pt, opacity=0.9] (1.72,3.22) rectangle (2.48,3.98);

		% ============================================================
		%  9. 细边框（墨线）
		% ============================================================
		\draw[fcInk, line width=0.7pt, opacity=0.45]
			([shift={(0.85,0.85)}]current page.south west) rectangle
			([shift={(-0.85,-0.85)}]current page.north east);

		% ============================================================
		% 10. 顶部眉题（课程英文小字 + 自适应墨线）
		% ============================================================
		\node[anchor=north, text=fcInk, align=center]
			(fcEyebrow) at ([yshift=-3.4cm]current page.north) {%
			{\sffamily\fontsize{12}{15}\selectfont \uppercase{\notecourseen}}
		};
		\draw[fcInk, line width=1.0pt] ($(fcEyebrow.west)+(-0.25cm,0)$) -- ++(-1.1cm,0);
		\draw[fcInk, line width=1.0pt] ($(fcEyebrow.east)+(0.25cm,0)$) -- ++(1.1cm,0);

		% ============================================================
		% 11. 主标题（居中，使用 \notename）
		% ============================================================
		\node[anchor=center, align=center, text=fcInk]
			at ([yshift=-8.3cm]current page.north) {%
			{\fontsize{38}{48}\sffamily\bfseries \notename}
		};

		% ============================================================
		% 12. 菱形 + 双细线
		% ============================================================
		\coordinate (C) at ([yshift=-11.6cm]current page.north);
		\draw[fcInk, line width=1.1pt] (C) ++(-2.2cm,0) -- ++(1.55cm,0);
		\draw[fcInk, line width=1.1pt] (C) ++(0.65cm,0) -- ++(1.55cm,0);
		\node[fill=fcInk, minimum size=0.24cm, inner sep=0pt, rotate=45] at (C) {};

		% ============================================================
		% 13. 副标题（使用 \notesubtitle）
		% ============================================================
		\node[anchor=north, text=fcInk, align=center]
			at ([yshift=-12.4cm]current page.north) {
			{\fontsize{15}{20}\sffamily \notesubtitle}
		};

		% ============================================================
		% 14. 底部作者、日期、附加信息
		% ============================================================
		\coordinate (B) at ([yshift=2.7cm]current page.south);
		\draw[fcInk, line width=0.8pt, opacity=0.8] (B) ++(-1.1cm,0.95cm) -- ++(2.2cm,0);
		\node[anchor=south, text=fcInk, align=center] at (B) {
			\begin{tabular}{c}
				{\fontsize{19}{24}\sffamily\bfseries \noteauthor} \\[0.35cm]
				{\large \sffamily \today}
				\ifstrempty{\noteextra}{}{\\[0.3cm]{\normalsize\sffamily \noteextra}}
			\end{tabular}
		};

		% ============================================================
		% 15. 右下角落款（画作信息）
		% ============================================================
		\node[anchor=south east, text=fcInk, opacity=0.65, font=\footnotesize\itshape]
			at ([shift={(-1.35,1.25)}]current page.south east)
			{Huang Gongwang, Dwelling in the Fuchun Mountains, 1350};

		\end{scope}
	\end{tikzpicture}
\end{titlepage}

%
%  配色：fc* 前缀（绢素 + 水墨五色）。本文件用 \providecolor 兜底，
%        单独复制到任何笔记本也能编译；想整体换调改这里即可。
%
%  注意：整幅画面绘于 \begin{scope}[shift={(current page.south west)}] 内，
%        所有坐标以「页面左下角为原点、单位 cm」书写；脱离该 scope
%        直接写 (x,y)，图形会落到页面外而不可见。
%
%  可用字段（在主文件导言区用 \newcommand 预定义即可覆盖缺省值）：
%      \notename     主标题（必填，主模板已定义）
%      \noteauthor   作者（必填，主模板已定义）
%      \notecourseen 眉题英文（缺省 Course Notes）
%      \notesubtitle 副标题（缺省「学习笔记」）
%      \noteextra    底部附加信息（缺省不显示）
% ============================================================

% ---- 《富春山居图》取色（绢素 + 水墨） ----
\providecolor{fcSilk}{RGB}{236,226,202}     % 绢底：米黄
\providecolor{fcSilkDark}{RGB}{214,202,174} % 绢底：暗部
\providecolor{fcInk3}{RGB}{178,176,168}     % 淡墨（远山）
\providecolor{fcInk2}{RGB}{116,114,108}     % 中墨（近坡）
\providecolor{fcInk}{RGB}{54,52,48}         % 浓墨（树、点、标题）
\providecolor{fcSeal}{RGB}{168,52,44}       % 朱红钤印

% ---- 缺省字段（主文件已定义的同名命令不会被覆盖） ----
\providecommand{\notecourseen}{Course Notes}
\providecommand{\notesubtitle}{学\ 习\ 笔\ 记}
\providecommand{\noteextra}{}

\begin{titlepage}
	\begin{tikzpicture}[remember picture, overlay]
		\begin{scope}[shift={(current page.south west)}]

		% ============================================================
		%  1. 绢底（米黄，下缘略深；带细微斑驳）
		% ============================================================
		\shade[top color=fcSilk, bottom color=fcSilkDark] (0,0) rectangle (21,29.7);
		\foreach \x/\y in {3.2/26.4, 8.6/24.2, 13.4/27.0, 17.8/25.2, 6.0/21.0,
		                   11.2/19.4, 19.0/21.8, 2.0/13.6, 15.2/8.2}{
			\fill[fcInk3, opacity=0.10] (\x,\y) ellipse[x radius=0.9, y radius=0.35];
		}

		% ============================================================
		%  2. 远山（三层，越远越淡；披麻皴短线）
		% ============================================================
		\fill[fcInk3, opacity=0.40]
			(0,20.0) -- (2.4,22.6) -- (4.0,21.2) -- (6.2,23.8) -- (8.6,21.6)
			-- (11.0,24.2) -- (13.4,21.4) -- (15.8,23.4) -- (18.2,21.0)
			-- (21,22.4) -- (21,17.4) -- (0,17.4) -- cycle;
		\fill[fcInk3, opacity=0.62]
			(0,17.8) -- (2.8,20.0) -- (4.6,18.6) -- (7.0,20.8) -- (9.4,18.4)
			-- (12.2,20.6) -- (15.0,18.2) -- (18.0,20.0) -- (21,18.2)
			-- (21,15.0) -- (0,15.0) -- cycle;
		\fill[fcInk2, opacity=0.45]
			(0,15.4) -- (2.2,17.0) -- (4.0,15.8) -- (6.4,17.4) -- (8.8,15.6)
			-- (11.6,17.2) -- (14.4,15.4) -- (17.6,16.8) -- (21,15.0)
			-- (21,12.6) -- (0,12.6) -- cycle;
		% 披麻皴：山体上的短竖笔
		\foreach \x/\y/\l in {3.0/19.4/1.5, 5.4/20.2/1.8, 7.8/19.6/1.6, 10.4/20.4/2.0,
		                      13.0/19.0/1.7, 16.0/19.4/1.5, 18.6/18.6/1.4}{
			\draw[fcInk2, opacity=0.35, line width=0.9pt]
				(\x,\y) .. controls (\x+0.18,\y-\l*0.45) and (\x-0.12,\y-\l*0.75) .. (\x+0.10,\y-\l);
		}

		% ============================================================
		%  3. 江水（留白 + 数笔淡墨波纹）
		% ============================================================
		\foreach \y in {12.0, 11.2, 10.4}{
			\draw[fcInk3, opacity=0.42, line width=0.8pt]
				plot[domain=1.2:19.6, samples=40, smooth]
					(\x, {\y + 0.13*sin(deg(2.0*\x + 3*\y))});
		}

		% ============================================================
		%  4. 一叶扁舟
		% ============================================================
		\fill[fcInk]
			(13.4,11.15) .. controls (14.0,10.90) and (15.0,10.90) .. (15.7,11.15)
			.. controls (15.2,11.42) and (13.9,11.42) .. (13.4,11.15) -- cycle;
		\draw[fcInk, line width=0.9pt] (14.6,11.15) -- (14.6,11.95);
		\fill[fcInk] (14.6,12.15) circle[radius=0.13];

		% ============================================================
		%  5. 近景坡石（左右两岸，中墨勾勒 + 皴笔）
		% ============================================================
		\fill[fcInk2, opacity=0.75]
			(0,0) -- (0,7.2) -- (1.6,8.6) -- (3.2,7.8) -- (4.8,9.4)
			-- (6.4,8.2) -- (8.0,9.0) -- (9.2,7.4) -- (10.2,5.6)
			-- (11.0,3.2) -- (11.6,0) -- cycle;
		\fill[fcInk2, opacity=0.70]
			(21,0) -- (21,6.0) -- (19.4,7.4) -- (17.8,6.6) -- (16.2,8.0)
			-- (14.6,6.8) -- (13.2,5.2) -- (12.4,2.6) -- (12.0,0) -- cycle;
		\draw[fcInk, opacity=0.55, line width=1.1pt]
			(0,7.2) -- (1.6,8.6) -- (3.2,7.8) -- (4.8,9.4) -- (6.4,8.2)
			-- (8.0,9.0) -- (9.2,7.4) -- (10.2,5.6) -- (11.0,3.2);
		\draw[fcInk, opacity=0.55, line width=1.1pt]
			(21,6.0) -- (19.4,7.4) -- (17.8,6.6) -- (16.2,8.0) -- (14.6,6.8)
			-- (13.2,5.2) -- (12.4,2.6);
		\foreach \x/\y/\l in {2.0/7.6/1.4, 4.2/8.2/1.6, 6.6/7.4/1.3, 8.6/7.8/1.2,
		                      15.0/6.2/1.3, 17.0/6.8/1.5, 19.0/6.2/1.2}{
			\draw[fcInk, opacity=0.30, line width=0.8pt] (\x,\y) -- (\x+0.12,\y-\l);
		}

		% ============================================================
		%  6. 疏落的秋树（竖笔为干，横点为叶）
		% ============================================================
		\foreach \x/\h in {1.4/3.2, 2.2/4.3, 3.1/3.6, 3.9/2.8, 15.6/3.0, 16.5/4.0, 17.3/3.2}{
			\draw[fcInk, line width=1.4pt] (\x,8.6) -- (\x,{8.6+\h});
			\foreach \dx/\dy in {-0.42/0.30, -0.16/0.52, 0.16/0.44, 0.42/0.24}{
				\fill[fcInk, opacity=0.75]
					({\x+\dx},{8.6+\h+\dy}) ellipse[x radius=0.20, y radius=0.10];
			}
		}

		% ============================================================
		%  7. 标题衬底（一层薄绢，保证文字可读）
		% ============================================================
		\fill[fcSilk, opacity=0.50]
			(2.6,16.9) -- (17.4,17.3) -- (18.0,22.6) -- (16.4,23.6) -- (2.2,22.9) -- cycle;

		% ============================================================
		%  8. 朱红钤印（右上、左下各一）
		% ============================================================
		\fill[fcSeal, opacity=0.85] (18.3,25.0) rectangle (19.7,26.4);
		\draw[fcSilk, line width=0.7pt, opacity=0.9] (18.55,25.25) rectangle (19.45,26.15);
		\fill[fcSeal, opacity=0.80] (1.5,3.0) rectangle (2.7,4.2);
		\draw[fcSilk, line width=0.6pt, opacity=0.9] (1.72,3.22) rectangle (2.48,3.98);

		% ============================================================
		%  9. 细边框（墨线）
		% ============================================================
		\draw[fcInk, line width=0.7pt, opacity=0.45]
			([shift={(0.85,0.85)}]current page.south west) rectangle
			([shift={(-0.85,-0.85)}]current page.north east);

		% ============================================================
		% 10. 顶部眉题（课程英文小字 + 自适应墨线）
		% ============================================================
		\node[anchor=north, text=fcInk, align=center]
			(fcEyebrow) at ([yshift=-3.4cm]current page.north) {%
			{\sffamily\fontsize{12}{15}\selectfont \uppercase{\notecourseen}}
		};
		\draw[fcInk, line width=1.0pt] ($(fcEyebrow.west)+(-0.25cm,0)$) -- ++(-1.1cm,0);
		\draw[fcInk, line width=1.0pt] ($(fcEyebrow.east)+(0.25cm,0)$) -- ++(1.1cm,0);

		% ============================================================
		% 11. 主标题（居中，使用 \notename）
		% ============================================================
		\node[anchor=center, align=center, text=fcInk]
			at ([yshift=-8.3cm]current page.north) {%
			{\fontsize{38}{48}\sffamily\bfseries \notename}
		};

		% ============================================================
		% 12. 菱形 + 双细线
		% ============================================================
		\coordinate (C) at ([yshift=-11.6cm]current page.north);
		\draw[fcInk, line width=1.1pt] (C) ++(-2.2cm,0) -- ++(1.55cm,0);
		\draw[fcInk, line width=1.1pt] (C) ++(0.65cm,0) -- ++(1.55cm,0);
		\node[fill=fcInk, minimum size=0.24cm, inner sep=0pt, rotate=45] at (C) {};

		% ============================================================
		% 13. 副标题（使用 \notesubtitle）
		% ============================================================
		\node[anchor=north, text=fcInk, align=center]
			at ([yshift=-12.4cm]current page.north) {
			{\fontsize{15}{20}\sffamily \notesubtitle}
		};

		% ============================================================
		% 14. 底部作者、日期、附加信息
		% ============================================================
		\coordinate (B) at ([yshift=2.7cm]current page.south);
		\draw[fcInk, line width=0.8pt, opacity=0.8] (B) ++(-1.1cm,0.95cm) -- ++(2.2cm,0);
		\node[anchor=south, text=fcInk, align=center] at (B) {
			\begin{tabular}{c}
				{\fontsize{19}{24}\sffamily\bfseries \noteauthor} \\[0.35cm]
				{\large \sffamily \today}
				\ifstrempty{\noteextra}{}{\\[0.3cm]{\normalsize\sffamily \noteextra}}
			\end{tabular}
		};

		% ============================================================
		% 15. 右下角落款（画作信息）
		% ============================================================
		\node[anchor=south east, text=fcInk, opacity=0.65, font=\footnotesize\itshape]
			at ([shift={(-1.35,1.25)}]current page.south east)
			{Huang Gongwang, Dwelling in the Fuchun Mountains, 1350};

		\end{scope}
	\end{tikzpicture}
\end{titlepage}
